% ============================================================================
%  fig_fsm.tex - FSM de control del juego (control_juego.v)
%  Pong en FPGA - Electrónica Digital 1 - UNAL
%  Juan David Ortiz Muñoz, David Eduardo Chivata Arias, Daniel Felipe Castro González
%  Compilar con: pdflatex fig_fsm.tex
% ============================================================================
\documentclass[border=6pt]{standalone}
\usepackage[utf8]{inputenc}
\usepackage[T1]{fontenc}
\usepackage{lmodern}
\usepackage{newunicodechar}
\newunicodechar{¬}{\ensuremath{\neg}}
\newunicodechar{∧}{\ensuremath{\wedge}}
\usepackage{tikz}
\usetikzlibrary{automata,positioning,arrows.meta,calc}
\begin{document}
\begin{tikzpicture}[
    font=\sffamily\footnotesize,
    >={Stealth[length=2.3mm]},
    every state/.style={draw, line width=0.7pt, minimum size=2.1cm, align=center, fill=violet!10},
    tr/.style={->, line width=0.6pt},
    lbl/.style={font=\sffamily\scriptsize, align=center, fill=white, inner sep=1.5pt},
]
\node[state, initial, initial text={INIT\_GLOBAL}] (esp) at (0, 0)  {ESPERA\\\texttt{000}};
\node[state]                (jue) at (6, 0)  {JUEGO\\\texttt{001}};
\node[state]                (pun) at (6, -5) {PUNTO\\\texttt{010}};
\node[state, fill=red!8]    (fin) at (0, -5) {FIN\\\texttt{100}};
\node[state, fill=orange!12] (prf) at (12.5, 3.6) {PAUSA\_RF\\\texttt{011}};
\node[state, fill=orange!12] (pma) at (12.5, -3.6) {PAUSA\_\\MANUAL\\\texttt{101}};

\draw[tr] (esp) -- node[lbl, above=1pt] {RF\_OK\_EFF ∧ START\_PULSE} (jue);
\draw[tr] (jue) to[bend left=18] node[lbl, sloped, above] {¬RF\_OK\_EFF \quad (prioridad 1)} (prf);
\draw[tr] (prf) to[bend left=18] node[lbl, sloped, below] {RF\_OK\_EFF ∧ START\_PULSE} (jue);
\draw[tr] (jue) to[bend left=18] node[lbl, sloped, above] {PAUSE\_PULSE \quad (2)} (pma);
\draw[tr] (pma) to[bend left=18] node[lbl, sloped, below, pos=0.35] {RF\_OK\_EFF ∧ PAUSE\_PULSE} (jue);
\draw[tr] (jue) -- node[lbl, left, pos=0.4] {GOAL\_J1 $\oplus$ GOAL\_J2\\(prioridad 3)} (pun);
\draw[tr] (pun) -- node[lbl, above] {WIN\_NEXT} (fin);
\draw[tr] (pun) to[bend left=10] node[lbl, sloped, above, pos=0.45] {¬WIN\_NEXT (saque)} (esp);

\draw[tr] (jue) edge[loop above] node[lbl, above] {sin eventos} (jue);
\draw[tr] (prf) edge[loop above] node[lbl, above] {en otro caso} (prf);
\draw[tr] (pma) edge[loop below] node[lbl, below] {en otro caso} (pma);
\draw[tr] (fin) edge[loop below] node[lbl, below] {solo sale con RESET (INIT\_GLOBAL)} (fin);
\draw[tr] (esp) edge[loop above] node[lbl, above] {en otro caso} (esp);

\node[draw, dashed, rounded corners, align=left, font=\sffamily\scriptsize, anchor=north west, text width=6.6cm, execute at begin node=\raggedright]
  at (5.6, -7.0) {Códigos \texttt{110} y \texttt{111} no usados: la lógica de siguiente estado los manda a ESPERA y \texttt{supervisor\_fallos} activa FAULT\_RECOVER.\\[2pt]
  En F5/F6 RF\_OK\_EFF pasa a ser la señal DETECT de los dos mandos cableados (PAUSA\_RF = mando desconectado).};
\end{tikzpicture}
\end{document}
